A speculative drafter that reads its target's hidden states becomes mismatched whenever that target is fine-tuned or replaced, and the usual remedy is to train a new drafter. We asked how much of that adaptation can be moved into the drafter's input interface instead. RelaySpec trains one linear map per target feature stream, keeps the draft transformer, fusion projection, normalization, embeddings and output head frozen, and folds the maps into the existing projection so that inference runs the original architecture with one updated matrix.

The evidence supports a specific and bounded claim. For a fine-tuned target the maps improve every adapter we tested at two model scales, and the gain tracks acceptance length rather than any change to the decoding path. For a replaced target within a model family, reconstruction of the original interface recovers essentially all of a purpose-built native drafter on mathematical reasoning and most of it elsewhere, without any target-specific drafter training. Across families and tokenizers the same construction still accelerates the target but no longer matches a native drafter, which we report as the boundary of the method rather than a gap to be explained away. The construction is not specific to DFlash: three maps on a frozen EAGLE-3 drafter reproduce both results, including the pairing of objective to setting.

Two findings explain why this works and are worth separating from the throughput numbers. The factorization into one map per stream is doing real work, since a single matrix of the same size, initialized to the same function, reaches far less under a matched objective. And the useful part of the learned correction is directional: the component perpendicular to the original fused context recovers nearly all of the benefit, while the parallel component leaves acceptance at the unadapted level, which is what RMSNorm predicts. The maps redirect a usable context rather than undoing the change the fine-tune introduced.

The practical consequence is that supporting a new target variant becomes a routine operation. Adaptation uses two orders of magnitude less data than the published recipe for training the drafter, is already effective from as few as 16 examples, adds 0.33 GiB of GPU memory, and leaves the draft transformer shared across targets. The limits we did not overcome are cross-tokenizer transfer, code and chat workloads under transfer, and high-concurrency serving where acceptance length stops being the bottleneck. Extending the interface to those settings, and to longer contexts and other accelerators, is the natural next step.
